\section{Pure type systems}

\Aclp{PTS}, originally introduced independently by Terlouw and Berardy and popularized by \citet{B91},
  provide a unifying framework for a number of popular type systems (and logics), including notably \coc and (the core) \mltt.
The idea is to parameterize the syntax and rules of a language by a notion of sorts (or universes).
Additionally, the different notions of terms and types (and kinds, and so on) are unified into a single notion of expressions.
We present here a version of \ac{PTS} extended with explicit substitutions and typed judgmental equality, following \citet{FP15}.

More formally, a \ac{PTS} is a typed $\lambda$-calculus parameterized by the following data:
\begin{enumerate}
\item A set $\St$ of \textit{sorts}, denoted as $\tmcst s$;
\item A binary relation $\Ax \subseteq \St \times \St$ of typing \textit{axioms}, denoted as $(\tmcst s_1 : \tmcst s_2)$;
\item A ternary relation $\Ru \subseteq \St \times \St \times \St$ of \textit{rules}, defining the allowed function spaces.
\end{enumerate}
We refer to the triple $(\St, \Ax, \Ru)$ as the \ac{PTS}'s \textit{specification}.
The most common example of \acp{PTS} is Barendregt's $\lambda$-cube, which consists of eight languages, ranging from \stlc to \coc.
These systems are defined by two sorts $*, \square$, a single axiom $* : \square$, and rules of the form $(\tmcst s_1, \tmcst s_2, \tmcst s_2)$, all including $(*, *, *)$.
Another major example is \mltt, obtained by taking $\St = \Nat$, axioms $i : j$ if $i \leq j$, and rules $(i, j, k)$ if $i, j \leq k$.
For the remainder of this section, we assume a fixed, arbitrary specification.


\subsection{Syntax}

We begin be looking at the syntactic side of \ac{PTS}.
Objects in a \ac{PTS} consists of the following syntactic categories:
\[
\begin{array}{llcrcl}
  \mbox{Expressions}       & \Exp & \ni & M,N,A,B,C      & ::=  & x_i \mid \tmcst s \mid \tppi A B \mid \tmlam M \mid \tmapp M N \mid \tmclo \sigma M \\
  \mbox{Contexts}          & \Ctx & \ni & \Gamma, \Delta & ::=  & \ctxempty \mid \ctxext\Gamma A \\
  \mbox{Substitutions}     & \Sub & \ni & \sigma, \tau   & ::=  & \subempty \mid \subid \mid \subwk \mid \subext \sigma M \mid \subclo \sigma \tau
\end{array}
\]
Although terms and types are unified into a single category of expressions, we distinguish them by writing $M,N$ for terms and $A,B,C$ for types.
The base cases of expressions are variables $x_i$, represented with de Bruijn indices, and sorts $\tmcst s \in \St$.
Function spaces $\tppi A B$ is the only type constructor, and can represent simple, dependent, or polymorphic functions, depending on the rule $r \in \Ru$ used in its construction.
The associated introduction form is the usual $\tmlam M$, which introduces a new variable allowed to occur in the body $M$ by shifting all de Bruijn indices up by one.
Likewise, the elimination form is the usual function application $\tmapp M N$ of a function $M$ to an argument $N$.
Last, $\tmclo \sigma M$ denotes the application of a substitution $\sigma$ to a term $M$.

Contexts are simply lists of types, growing to the right.
We write $\ctxempty$ for the empty context and $\ctxext \Gamma A$ for the extension of context $\Gamma$ with a variable of type $A$.
The de Bruijn indices of variables correspond to their relative positions in the context.

The empty substitution $\subempty$ allows moving a closed expression to an arbitrary context.
The identity substitution $\subid$ intuitively does nothing, and is mainly useful for formulating inference rules.
The weakening substitution $\subwk$ shifts all de Bruijn indices up by one.
The substitution extension $\subext \sigma M$ extends the domain of a substitution $\sigma$ from $\Gamma$ to $\ctxext \Gamma A$, provided that $M$ has type $A$.
Finally, substitutions can be compose with the form $\subclo \sigma \tau$, where $\tau$ is applied first and then $\sigma$.





\subsection{Judgments}

\begin{figure}
  \raggedright
  \noindent
  $\boxed{\wfexp \Gamma M A}$ -- Expression $M$ has type $A$ in context $\Gamma$.
  \[
  \begin{array}{c}
    \infer{
      \wfexp \Gamma {x_i} A
    }{
      \ctxlookup \Gamma {x_i} A
    }
    \qquad
    \infer{
      \wfexp \Gamma {\tmcst s_1} {\tmcst s_2}
    }{
      \wfctx \Gamma
      &
      (\tmcst s_1 : \tmcst s_2) \in \Ax
    }
    \qquad
    \infer{
      \wfexp \Gamma {\tmclo \sigma M} {\tmclo \sigma A}
    }{
      \wfexp \Delta M A
      &
      \wfsub \Gamma \sigma \Delta
    }
    \\[8pt]
    \infer{
      \wfexp \Gamma {\tppi A B} {\tmcst s_3}
    }{
      \wfexp \Gamma A {\tmcst s_1}
      &
      \wfexp {\ctxext \Gamma A} B {\tmcst s_2}
      &
      r = (\tmcst s_1, \tmcst s_2, \tmcst s_3) \in \Ru
    }
    \qquad
    \infer{
      \wfexp \Gamma {\tmlam M} {\tppi A B}
    }{
      \wftyp \Gamma \tppi A B
      &
      \wfexp {\ctxext \Gamma A} M B
    }
    \\[8pt]
    \infer{
      \wfexp \Gamma {\tmapp M N} {\tmclo {\subext \subid N} B}
    }{
      \wftyp \Gamma {\tppi A B}
      &
      \wfexp \Gamma M {\tppi A B}
      &
      \wfexp \Gamma N A
    }
    \qquad
    \infer{
      \wfexp \Gamma M B
    }{
      \wfexp \Gamma M A
      &
      \eqtyp \Gamma A B
    }
  \end{array}
  \]
  %
  \noindent
  $\boxed{\wftyp \Gamma A}$ -- $A$ is a well-formed type in context $\Gamma$.
  \[
  \begin{array}{c}
    \infer{
      \wftyp \Gamma {\tmcst s}
    }{
      \wfctx \Gamma
      &
      \tmcst s \in \St
    }
    \qquad
    \infer{
      \wftyp \Gamma A
    }{
      \wfexp \Gamma A {\tmcst s}
      &
      \tmcst s \in \St
    }
    \qquad
    \infer{
      \wftyp \Gamma {\tmclo \sigma {\tmcst s}}
    }{
      \wfsub \Gamma \sigma \Delta
      &
      \tmcst s \in \St
    }
  \end{array}
  \]
  %
  \noindent
  $\boxed{\wfsub \Gamma \sigma \Delta}$ -- $\sigma$ is a well-formed substitution from $\Delta$ to $\Gamma$.
  \[
  \begin{array}{c}
    \infer{
      \wfsub \Gamma \subempty \ctxempty
    }{
      \wfctx \Gamma
    }
    \qquad
    \infer{
      \wfsub \Gamma \subid \Gamma
    }{
      \wfctx \Gamma
    }
    \qquad
    \infer{
      \wfsub {\ctxext \Gamma A} \subwk \Gamma
    }{
      \wfctx {\ctxext \Gamma A}
    }
    \qquad
    \infer{
      \wfsub \Gamma {\subext \sigma M} {\ctxext \Delta A}
    }{
      \wfsub \Gamma \sigma \Delta
      &
      \wfexp \Delta A {\tmcst s}
      &
      \wfexp \Gamma M {\tmclo \sigma A}
    }
    \\[8pt]
    \infer{
      \wfsub {\Gamma_1} {\subclo \sigma \tau} {\Gamma_3}
    }{
      \wfsub {\Gamma_1} \sigma {\Gamma_2}
      &
      \wfsub {\Gamma_2} \tau {\Gamma_3}
    }
    \qquad
    \infer{
      \wfsub \Gamma \sigma {\Delta'}
    }{
      \wfsub \Gamma \sigma \Delta
      &
      \eqctx \Delta {\Delta'}
    }
  \end{array}
  \]
  %
  \noindent
  $\boxed{\eqexp \Gamma M N A}$ -- $M$ and $N$ are equivalent expressions of type $A$ in context $\Gamma$.
  \[
  \begin{array}{c}
    \infer{
      \eqexp \Gamma {\tmapp {(\tmlam M)} N} {\tmclo {\subext \subid N} M} {\tmclo {\subext \subid N} B}
    }{
      \wfexp {\ctxext \Gamma A} M B
      &
      \wfexp \Gamma N A
    }
    \qquad
    \infer{
      \eqexp \Gamma M {\tmlam {(\tmapp {\tmclo \subwk M} {x_0})}} {\tppi A B}
    }{
      \wfexp \Gamma M {\tppi A B}
    }
  \end{array}
  \]
  %
  $\boxed{\eqtyp \Gamma A B}$ -- $A$ and $B$ are equal types in context $\Gamma$.
  \[
  \begin{array}{c}
    \infer{
      \eqtyp \Gamma {\tmcst s} {\tmcst s}
    }{
      \wfctx \Gamma
      &
      \tmcst s \in \St
    }
    \qquad
    \infer{
      \eqtyp \Gamma A B
    }{
      \eqexp \Gamma A B {\tmcst s}
      &
      \tmcst s \in \St
    }
    \qquad
    \infer{
      \eqtyp \Gamma {\tmclo \sigma {\tmcst s}} {\tmcst s}
    }{
      \wfsub \Gamma \sigma \Delta
      &
      \tmcst s \in \St
    }
  \end{array}
  \]
  \caption{Selected rules of \ac{PTS} judgments}
  \label{fig:pts-judg}
\end{figure}


Next, we discuss the judgments of \ac{PTS}.
Selected rules are shown in \Cref{fig:pts-judg}, a complete definition is given in \Cref{app:pts-syntax}.

Each syntactic category has two associated judgments: well-formedness (or typing) and equality judgments.
We use typed equivalence, similar to \citet{A06}.
We write $\wfexp \Gamma M A$ to mean that $M$ is a well-formed term of type $A$ in context $\Gamma$,
  and $\eqexp \Gamma M {M'} A$ to mean that $M$ and $M'$ are equal terms of type $A$ in context $\Gamma$.
For substitution, the judgment $\wfsub \Gamma \sigma \Delta$ states that $\sigma$ is a well-formed substitution with domain $\Delta$ and co-domain $\Gamma$.
The context formation judgment $\wfctx \Gamma$ is slightly different in that it simply states that $\Gamma$ is well-formed, but does have a component analogous to a type.

In addition the the well-formedness and equality judgments for the three syntactic categories, we define similar judgments for types.
The judgment $\wftyp \Gamma A$ states that $A$ is a well-formed type, which means that it either has a sort, or is itself a sort, and similarly for the type equality $\eqtyp \Gamma A B$.
We need those judgments because there can be sorts that do not themselves have a sort.
For example, in the systems of the $\lambda$-cube, the sort $\square$ is a well-formed sort that does not itself have a sort.
With these judgments, we take a step away from the judgmental equality proposed by \citet{A06}, where no separate notion of type equality is considered.
This is necessary because we use explicit substitutions, so we need a way to validate that $\tmclo \sigma {\tmcst s} \equiv \tmcst s$ even when $\tmcst s$ does not have a sort.

Last, we use a context lookup judgment, denoted $\ctxlookup \Gamma {x_i} A$.

The characteristic rules of \acp{PTS} are those parameterized by the specification:
\[
\infer{
  \wfctx {\ctxext \Gamma A}
}{
  \wfctx \Gamma
  &
  \wfexp \Gamma A {\tmcst s}
  &
  \tmcst s \in \St
}
\qquad
\infer{
  \wfexp \Gamma {\tmcst s_1} {\tmcst s_2}
}{
  (\tmcst s_1 : \tmcst s_2) \in \Ax
}
\qquad
\infer{
  \wfexp \Gamma {\tppi A B} {\tmcst s_3}
}{
  \wfexp \Gamma A {\tmcst s_1}
  &
  \wfexp {\ctxext \Gamma A} B {\tmcst s_2}
  &
  r = (\tmcst s_1, \tmcst s_2, \tmcst s_3) \in \Ru
}
\]
In the \ac{PTS} setting, ``types'' consists of sorts and any expression that has a sort.
The above rules specify precisely those expressions that can have a sort and ensure that contexts really do consist exclusively of well-sorted types.

Typing rules for functions and function application are typical of systems based on natural deductions.
We take extra care to ensure that types are well-formed via the additional premise that the function space is a well-formed type.
These premises are needed to ensure that the presupposition lemma (see \Cref{lemma:pts-presupposition}) holds.

The notion of equality for expressions stems from $\beta$-reductions and $\eta$-expansions.
Because we use explicit substitutions, the equality judgment must also specify how substitutions are propagated into expressions.
In addition, equality includes the usual congruence rules, as well as rules stating reflexivity, symmetry, and transitivity, i.e. that equality is an equivalence relation.

For equality of types, the main challenge is propagating substitutions in sorts.
Since sorts contain no variables, we really only need to forget about said substitutions.






\subsection{Syntactic properties}

A key advantage of \acp{PTS} is that the general framework allows us to prove some desirable properties parametrically in the specification.
In this way, we obtain generic proofs that can later be instantiated to concrete languages, thus reducing the proof burden of developping languages.
For an arbitrary \ac{PTS}, we can only prove simple properties, such as the Church-Rosser theorem, substitution properties, and subject reduction.
For our purposes, the main lemma is presupposition:
\begin{lemma}[Presupposition]
  \label{lemma:pts-presupposition}
  Every components of a judgment must be well-formed.
  Specifically:
  \begin{enumerate}
  \item If $\wfexp \Gamma M A$, then $\wfctx \Gamma$ and $\wftyp \Gamma A$.
  \item If $\eqexp \Gamma M {M'} A$, then $\wfctx \Gamma$, $\wftyp \Gamma A$, $\wfexp \Gamma M A$, and $\wfexp \Gamma {M'} A$.
  \end{enumerate}
  And similarly for other judgments.
  Moreover, the size of the resulting judgments is bounded by the size of the input judgments.
  \AG{This last part about size is true, but maybe not needed.  It would also be a pain to mechanize.}
\end{lemma}
\begin{proof}
  By induction on the input derivation.
  Most cases are straightforward.
  In the axiom case, we have a proof of $\wfexp \Gamma {\tmcst s_1} {\tmcst s_2}$,
    and we need the type well-formedness judgment to infer $\wftyp \Gamma {\tmcst s_2}$, as we cannot ensure the existence of an axiom $\tmcst s_2 : \tmcst s_3$.
\end{proof}

More interesting properties, like normalization, cannot be proven generically because they do not hold for all \ac{PTS}, like Girard's System U or Martin-L{\"o}f's \texttt{Type} : \texttt{Type} system.
However, we can prove such properties based on some conditions.
This is the approach taken by \citet{MW96}, who give a reducibility argument based on the existence of certain sets,
  and by \citet{FP15}, who give a \ac{NbE} argument for an arbitrary predicative \ac{PTS}.
